\documentclass[journal]{IEEEtran}
\usepackage{amsmath,amssymb}
\usepackage{graphicx}
\usepackage{booktabs}
\usepackage{multirow}
\usepackage{algorithm}
\usepackage{algpseudocode}
\usepackage{url}
\usepackage{cite}

\begin{document}

\title{Detection Is Not Investigation: Why RL-as-Classifier Fails in
Network Intrusion Detection, and What Separating the Two Does and Does
Not Buy}

\author{\IEEEauthorblockN{Author Name}
\IEEEauthorblockA{Affiliation withheld for review}
\thanks{Manuscript prepared for IEEE Transactions submission.}}

\markboth{IEEE Transactions on ...}{}

\maketitle

\begin{abstract}
A growing body of work asks a reinforcement-learning (RL) policy to own the
benign/attack decision itself. We screened ten such graph/RL
intrusion-detection designs and reproduce the failure: after correcting a
reward-scalarization bug, the best candidate still trailed the strongest
gradient-boosted-tree baseline by 20--53 macro-F1 points, and a
static-threshold classifier beat it on most (dataset, seed) pairs. We argue
the cause is structural rather than a tuning accident, and build the
redesign that diagnosis implies. In AegisGraph-X a nine-expert pool supplies
detection, a graph mixture-of-experts router combines it, and a budgeted RL
layer decides only whether to commit, acquire telemetry, escalate, isolate,
or defer, never re-deriving the label.

The redesign removes the failure mode but validates almost none of its own
machinery, and that is the contribution. Over ten seeds at a pinned training
budget the backbone reaches macro-F1 $0.940$, $0.997$, and $0.949$.
Equivalence testing establishes agreement with the strongest tree baseline
to within $0.001$ macro-F1 on two corpora, but fails on CICIoT2023, where
the backbone is reliably worse by $1.2$ points. The added machinery then
fails its own ablation: the single best expert \emph{significantly
outperforms} the full nine-expert pool ($p = 0.008$), the learned router is
indistinguishable from an unweighted average ($p = 0.90$), the RL layer
never spends its telemetry budget, conformal FPR control moves alert volume
by $0.2\%$, and the student's hard-class components cost $5.6$ points
against plain distillation. Only calibration survives its own ablation. A
closed-form decision-theoretic policy computed from the same calibrated
posterior and the same reward table, with no training, never isolates on
any dataset and significantly outperforms the trained backbone on two of
three -- so the RL layer's worst failure is attributable to training, not to
the reward specification it was trained against.
\end{abstract}

\begin{IEEEkeywords}
Intrusion detection, mixture of experts, gradient-boosted trees, knowledge
distillation, reinforcement learning, sequential decision-making, calibration,
conformal prediction, leakage-safe evaluation.
\end{IEEEkeywords}

\section{Introduction}
\label{sec:intro}

Network intrusion detection systems (IDS) are increasingly asked to do more
than emit a single benign/attack label: a security operations center (SOC)
analyst must decide, flow by flow and host by host, whether to gather more
telemetry, escalate to a human, isolate a host, or let a decision wait --
each option carrying a different cost and a different risk of missing or
falsely flagging an attack. Static classifiers, however well they score on a
held-out test split, are silent on this sequential, budget-constrained
investigation process. At the same time, a large and growing body of work has
tried to make reinforcement learning (RL) \emph{itself} the classifier --
letting an RL policy learn, from scratch, the benign/attack decision boundary
that decades of supervised tree-ensemble and deep-learning research have
already solved to a high standard. Our own earlier screening of ten candidate
graph/RL intrusion-detection designs (Section~\ref{sec:method_history})
reproduces this failure directly: even after diagnosing and correcting a
genuine reward-tuning pathology, the best RL-as-classifier candidate we could
validate still trailed the strongest gradient-boosted-tree baseline by
20--53 macro-F1 points per dataset, and its own ablation showed a trivial
static-threshold classifier beat it on the majority of (dataset, seed)
combinations. The root cause is structural, not a tuning accident: a policy
trained on a few hundred offline episodes or a few dozen on-policy updates
cannot out-learn a 300-tree gradient-boosted ensemble at the same supervised
task, and asking it to try wastes RL's actual comparative advantage --
sequential, budget-aware decision-making -- on a problem it is not suited to
win.

\textbf{This paper's premise.} Reinforcement learning should not replace a
strong classifier; it should sit on top of one. We propose
\textbf{AegisGraph-X}, a cost-aware intrusion-detection architecture that
separates \emph{what} to decide from \emph{when and how thoroughly} to
decide it: a high-performance pool of tree, boosting, tabular-neural,
temporal, and graph classifiers (Section~\ref{sec:aegis_experts}) supplies
the classification signal, a mixture-of-experts (MoE) router
(Section~\ref{sec:aegis_router}) combines them per-instance, a tree-guided
distillation stage (Section~\ref{sec:aegis_distill}) produces a compact
deployable student, and a budgeted graph-RL investigation layer
(Section~\ref{sec:aegis_rl}) decides only whether to commit to the
classifier's calibrated prediction, request more telemetry, expand the graph
neighborhood, escalate to a human analyst, isolate a host, adjust its
operating threshold, or defer -- never re-deriving the class label itself.
This division of labor is enforced structurally, not just by convention: the
RL policy's only classification-relevant action (COMMIT) reads its label
directly from the pretrained, calibrated ensemble, so the achievable
classification quality is bounded below by the classifier's own quality
regardless of how well or poorly the RL policy trains
(Section~\ref{sec:aegis_rl}).

\textbf{Research gap.} We identify three gaps that this paper addresses
directly: (1) a large share of the graph/RL-for-IDS literature evaluates an
RL policy as if it were a competitive classifier, without comparing it
against the strong tree/boosting baselines that already solve most public IDS
benchmarks near ceiling; (2) when RL genuinely adds sequential-decision value
(delay, telemetry cost, false-alarm burden), that value is rarely separated
from -- and is often conflated with -- classification accuracy, making it
impossible to tell whether a reported gain is a real operational improvement
or an artifact of a favorable reward design; (3) reward-scale comparisons
across RL methods are rarely made fair: a method that tunes its own reward
weights on validation data is not directly comparable, on raw cumulative
reward, to methods evaluated under different or default weights.

\textbf{Contributions.}
\begin{itemize}
\item A structural fix -- not merely a tuning fix -- to the failure mode
identified in our own prior screening of ten graph/RL designs: AegisGraph-X's
COMMIT action always reads its label from a separately-trained, calibrated
expert ensemble, so the RL layer cannot degrade classification quality by
construction, only decision-cost and investigation behavior
(Section~\ref{sec:aegis_rl}).
\item A five-layer cost-aware architecture -- expert pool, graph
mixture-of-experts router, tree-guided neural distillation, budgeted graph-RL
investigation, and calibration/low-FPR control -- built and evaluated under
the same leakage-safe dataset-screening protocol as our prior study
(Section~\ref{sec:datasets}, reused unchanged), with every component
individually ablated (Section~\ref{sec:results}).
\item A validation-only, three-way comparison (backbone-only vs. +RL
investigation vs. +conformal safety shield) that selects the deployed
configuration on validation evidence alone, and reports honestly when the RL
or safety-shield layer improves decision utility but not classification
accuracy, or vice versa (Section~\ref{sec:results}).
\item A dual-mode reinforcement-learning evaluation (common fixed reward
weights vs. validation-tuned reward weights, reported separately) that
prevents the raw-cumulative-reward comparability error our prior manuscript
disclosed as a limitation (Section~\ref{sec:results}).
\item A repeated-seed comparison against 41 other classical, deep, graph, and
RL/decision models (37 baselines, our own prior offline-safe-RL design kept
as baseline \#39 rather than discarded, and 3 proposed-variant ablation
points), reporting exact win/tie/loss counts with paired significance
testing, calibration quality, and runtime/memory
measurements (Section~\ref{sec:results}).
\end{itemize}

\section{Related Work}
\label{sec:related}

\textbf{IDS benchmarks.} UNSW-NB15~\cite{moustafa2015unswnb15},
CIC-IDS-2017~\cite{sharafaldin2018cicids}, ToN-IoT~\cite{alsaedi2020toniot},
and CICIoT2023~\cite{neto2023ciciot2023} are among the most widely used
public IDS corpora. Several are near-saturated for binary classification
under strong tree ensembles, which motivates multiclass and minority-recall
reporting alongside binary accuracy -- and which, as
Section~\ref{sec:limitations} discusses, materially limits what a binary
comparison at the ceiling can establish.

\textbf{Tree/boosting IDS.} Gradient-boosted and bagged tree ensembles --
XGBoost~\cite{chen2016xgboost}, LightGBM~\cite{ke2017lightgbm},
CatBoost~\cite{prokhorenkova2018catboost}, random
forests~\cite{breiman2001randomforests} -- remain the strongest published
classifiers on most public IDS benchmarks, frequently reaching near-ceiling
binary macro-F1. This mirrors the broader finding that tree ensembles remain
competitive with, and often superior to, deep models on tabular
data~\cite{grinsztajn2022trees,shwartzziv2022tabular}. We treat this as a
fact to build on rather than a target to dethrone with RL, using six such
ensembles as the core expert pool (Section~\ref{sec:aegis_experts}).

\textbf{Deep tabular IDS.} MLP, CNN/RNN-family, Transformer, and
tabular-attention models -- FT-Transformer~\cite{gorishniy2021revisiting},
TabTransformer~\cite{huang2020tabtransformer}, SAINT~\cite{somepalli2021saint}
-- close some of the gap to tree ensembles on selected corpora at higher
training cost. We include an FT-Transformer-style tabular expert and a
temporal convolutional expert~\cite{bai2018tcn} in the pool to capture this
family's signal.

\textbf{Graph IDS.} GNN-based IDS work encodes flows as host or protocol-role
graphs and applies GCN~\cite{kipf2017gcn},
GraphSAGE~\cite{hamilton2017graphsage}, or GAT~\cite{velickovic2018gat}
variants for node- or edge-level classification;
E-GraphSAGE~\cite{lo2022egraphsage} is representative of the flow-based
edge-feature line. This work rarely evaluates a decision cost beyond the
classification head. AegisGraph-X includes a GraphSAGE-style protocol-role
graph expert and reuses the graph construction from our prior study
unchanged (Section~\ref{sec:datasets}).

\textbf{Reinforcement learning for cybersecurity.} RL has been applied to
intrusion detection largely by recasting per-flow classification as a
single-step ``bandit'' decision or a multi-step investigation task in which
the policy itself must learn the class boundary from
scratch~\cite{nguyen2019drlcyber}; our prior ten-candidate screening
(Section~\ref{sec:method_history}) confirms this consistently underperforms
mature supervised classifiers at the seed/update budgets tractable for a
single study. We instead keep RL out of the classification decision
entirely, using it only for investigation-order, telemetry acquisition, and
escalation/isolation policy, with PPO~\cite{schulman2017ppo} as the trainer.

\textbf{Mixture of experts and stacking.} Per-instance gating over a pool of
specialized sub-models dates to adaptive mixtures of local
experts~\cite{jacobs1991moe} and is now standard in large-scale language
modeling~\cite{shazeer2017moe,fedus2022switch}, but is under-explored for
tabular/graph intrusion detection, where ``experts'' are typically an ad hoc
average rather than a learned, context-conditioned router. The closer
antecedent for what our router actually does is stacked
generalization~\cite{wolpert1992stacked}: fitting a combiner on held-out
base-learner predictions. That literature also anticipates our negative
result -- learned combiners frequently fail to beat simple averaging at
modest sample sizes, which is the regime our router occupies
(Section~\ref{sec:results}).

\textbf{Calibration and safe decision-making.} Temperature
scaling~\cite{guo2017calibration}, isotonic
regression~\cite{zadrozny2002isotonic}, and split-conformal
prediction~\cite{angelopoulos2021conformal} give calibrated or
distribution-free signals usable for false-positive-rate control. We use
these to calibrate the ensemble probability before any threshold is applied
(Section~\ref{sec:aegis_calib}), and use a conformal prediction-set-size
signal as the basis of Candidate C's safety shield
(Section~\ref{sec:aegis_rl}).

\textbf{Sequential acquisition and the value of information.} Deciding what
to observe before committing to a decision is a classical problem: Wald's
sequential probability ratio test~\cite{wald1945sequential} is the canonical
single-hypothesis case, and information value theory~\cite{howard1966voi}
gives the general rule -- acquire evidence only when its expected reduction
in Bayes risk exceeds its cost. We adapt budgeted acquisition to
graph-structured telemetry as one of the RL layer's action choices.
Section~\ref{sec:limitations} returns to the question this framing raises and
that our experiments do not settle: whether a learned policy is needed here
at all, given that a calibrated posterior plus an explicit cost table admits
an analytic solution.

\textbf{Offline/safe RL.} Conservative Q-learning~\cite{kumar2020cql} and
related behavior-constrained offline methods~\cite{fujimoto2019bcq} avoid
unsafe online exploration by learning from logged interaction data; our prior
offline-safe-RL design (Section~\ref{sec:method_history}) is retained as
baseline \#39, a direct comparison point for the on-policy budgeted
investigation layer.

\textbf{Distillation.} The compact student follows standard knowledge
distillation~\cite{hinton2015distilling}, with focal loss~\cite{lin2017focal}
and supervised contrastive learning~\cite{khosla2020supcon} as the
class-imbalance components ablated in Section~\ref{sec:results}.

\section{Dataset Screening and POMDP Construction}
\label{sec:datasets}

\subsection{Fourteen-dataset screening}
Table~\ref{tab:dataset_screening} summarizes the full recursive inventory of
the fourteen nominal dataset folders made available for this study, which
resolve to twelve unique dataset families (two folder pairs are duplicate
copies of the same underlying file). Each dataset was screened against seven
predeclared criteria: (1) leakage-safe split feasibility, (2) usable
benign/attack labels, (3) graph/temporal/sequential construction feasibility,
(4) absence of severe saturation/duplication risk, (5) meaningful binary
and/or multiclass evaluation, (6) sufficient samples for repeated-seed
testing, and (7) support for a meaningful sequential decision process. Full
screening detail, including two datasets initially misjudged from a
non-representative row sample and corrected by a full-column rescan, is
reported in the supplementary \texttt{dataset\_selection\_report.md}.

\input{../tables/table_I_dataset_screening.tex}

\subsection{Final dataset selection}
CICIoT2023 (modern large-scale IoT/IIoT, 34 attack/benign classes, native
curator-provided train/validation/test split, respected here with only
post-hoc exact-duplicate removal across splits), ToN-IoT-Network
(IoT/IIoT telemetry-fusion testbed retaining real source/destination IPs,
enabling genuine host-flow graphs and host-isolation semantics), and
UNSW-NB15 (the most widely benchmarked general network-IDS corpus among the
twelve, keeping repeated-seed and multi-baseline experiments tractable) are
retained as primary evidence. RT-IoT2022~\cite{sharmila2023rtiot2022} is
retained only as a supplementary cross-dataset-transfer check. Per explicit design decision, NSL-KDD and
APA-DDoS are used only as sanity-control datasets and never as primary
evidence, because of NSL-KDD's well-documented saturation/duplication history
and APA-DDoS's narrow single-attack-family, 151K-row scope. CIC-VPN2016 and
ISCX-VPN2016 were excluded outright: their labels encode traffic/application
type (chat, VoIP, streaming, VPN-vs-plain), not benign/attack, so they fail
criterion (2) for this task. ISCX-IDS-2012 was excluded for a 60.3\% duplicate
rate in a 20K-row sample, a serious leakage risk for this legacy corpus.
Bot-IoT was excluded purely on computational-feasibility grounds (15\,GB
across 74 shard files, with an attack-concentrated shard layout that makes a
trustworthy minority-class estimate require a full rescan outside this
study's compute budget) -- a scope limitation, not a claim of unsuitability.

\input{../tables/table_II_dataset_stats.tex}

\subsection{Leakage-safe preprocessing}
For every retained dataset, the split is fixed \emph{before} any
preprocessing is fit: session/flow-5-tuple-disjoint grouping is used when a
real session identifier is available (ToN-IoT-Network, RT-IoT2022, UNSW-NB15
falls back honestly to stratified-random since its release CSV carries only a
row index, not a real session key); scalers, encoders, and imputers are fit
on the training partition only; exact cross-split duplicates are removed from
validation/test (never from train); and CICIoT2023's curator-provided
train/validation/test file split is respected rather than reshuffled. Every
split manifest (row counts, class balance per split, deduplication counts) is
saved under \texttt{manifests/} for audit.

\subsection{Graph construction}
Two graph types are used depending on the presence of raw host identifiers:
a \emph{host-flow graph} (nodes = hosts, edges = flows) for ToN-IoT-Network,
and a \emph{protocol-role graph} (nodes = (protocol, port-group, role)
tuples) for CICIoT2023, UNSW-NB15, and RT-IoT2022, none of which release raw
IP addresses. Graphs are built over fixed-size windows of consecutive flows
(64 flows per window in all reported experiments); for graph baselines, node
features are the mean of incident flows' scaled feature vectors rather than a
bare degree count, giving the GNN baselines a fair, information-rich input.

\subsection{POMDP formulation (legacy environment)}
Table~\ref{tab:pomdp} summarizes the state, action, and reward structure of
the original 11-action POMDP, reused unchanged for the ten legacy candidates
and the RL/decision baselines (Section~\ref{sec:setup}); AegisGraph-X's own
redesigned 10-action environment is described separately in
Section~\ref{sec:aegis_rl}. An episode corresponds to one graph window: the
agent must dispose of every flow
in the window (classify, abstain, or escalate) while optionally spending a
bounded per-episode telemetry budget acquiring feature groups, node-level
statistics, or an expanded subgraph neighborhood first. Isolating a host has
a genuine consequence within the episode: subsequent flows from that host are
treated as contained (crediting prevented harm if the host was truly
malicious, penalizing disruption if it was not), which is what makes
isolation a real sequential decision rather than a relabeled classification.

\input{../tables/table_III_pomdp.tex}

\section{Proposed AegisGraph-X}
\label{sec:method}

\subsection{Why RL-as-classifier fails: summary of prior screening}
\label{sec:method_history}
Before designing AegisGraph-X, we screened ten candidate graph/RL designs in
which an RL policy directly owned the benign/attack classification decision
(11-action POMDP with independent CLASSIFY\_BENIGN/CLASSIFY\_ATTACK actions,
sharing a common one-layer GraphSAGE-style graph encoder and PPO-lite/offline
-CQL trainer so differences reflect each candidate's distinguishing module,
not unequal engineering effort). Validation-only screening selected an
offline conservative-Q-learning design (Candidate 3) -- \emph{not} the
task's a-priori top-priority candidate, evidence the selection was not
cherry-picked. Even after diagnosing and correcting a genuine reward-scalarization
failure (a composite validation objective's eight uncorrected terms let a
majority-class-collapsing policy score well; fixed with a continuous
minority-recall gate and median-of-3-training-seed trial aggregation -- full
diagnosis retained in \texttt{manifests/technical\_audit\_report.md} for the
record), Candidate 3's corrected test macro-F1 ($0.465$--$0.679$) still
trailed the strongest gradient-boosted-tree baseline by 20--53 points per
dataset, and its own ablation showed a trivial static-threshold classifier
beat the ``full'' policy on 7 of 9 (dataset, seed) combinations. The root
cause is structural: at the seed/update budgets tractable for one study, an
RL policy cannot out-learn a mature supervised ensemble at the classification
task itself. \textbf{We therefore do not retain Candidate 3 (or any
RL-as-classifier design) as the headline method}; it is kept only as baseline
\#39 (a direct, transparent comparison point) in every table below.

\subsection{Architecture overview}
AegisGraph-X separates classification from investigation into five layers.
\textbf{(1) Expert pool} (Section~\ref{sec:aegis_experts}): six tree/boosting
models, one tabular-neural model, one temporal model, and one graph model,
each independently trained to classify a flow. \textbf{(2) Graph
mixture-of-experts router} (Section~\ref{sec:aegis_router}): a learned gate
that mixes the pool's per-instance predictions, conditioned on graph/topology
context. \textbf{(3) Tree-guided neural distillation} (Section~\ref{sec:aegis_distill}):
compresses the router-mixed ensemble into one compact feed-forward student
for deployment settings where running all nine experts is too costly.
\textbf{(4) Budgeted graph-RL investigation} (Section~\ref{sec:aegis_rl}): a
PPO policy that decides whether to commit to the calibrated ensemble's
label, acquire more information, escalate, isolate, or defer -- never to
independently classify. \textbf{(5) Calibration and low-FPR control}
(Section~\ref{sec:aegis_calib}): temperature/isotonic/Platt scaling and
split-conformal thresholding, fit on validation data only, applied before any
operating threshold (including the RL layer's COMMIT decision) is used.
Class-imbalance handling (class-balanced focal loss, supervised contrastive
learning, loss-weighted hard-example resampling) is applied at the
distillation stage (Layer 3), where a compact student is most prone to
majority-class shortcut solutions; each technique is individually ablated
(Section~\ref{sec:results}).

\subsection{Layer 1: expert pool}
\label{sec:aegis_experts}
The pool comprises XGBoost, LightGBM, CatBoost, Random Forest, Extra Trees,
and HistGradientBoosting (six tree/boosting experts); one tabular-attention
expert (a lightweight FT-Transformer-style encoder over the flow's own
feature vector); one temporal expert (a temporal convolutional network over
a 16-flow sliding window); and one graph expert (a GraphSAGE-style encoder
over the same protocol-role/host-flow windowed graphs used elsewhere in this
study, Section~\ref{sec:datasets}, read out per-edge/per-flow). All nine are
trained with class-balanced weighting appropriate to their family (class
weights for trees/FT-Transformer/TCN; the same for the graph expert's
cross-entropy readout).

\textbf{Leakage-safe stacking.} Training data is split once into
\emph{expert\_fit} (70\%) and \emph{router\_fit} (30\%), contiguously rather
than by shuffled/stratified sampling -- a deliberate choice: the graph
expert's windows are built over runs of consecutive rows, so a shuffled
split would silently break window structure for that one expert while
leaving the other eight row-independent; a single contiguous split keeps the
procedure uniform across all nine. Experts are fit on \emph{expert\_fit} and
scored \emph{out-of-sample} on \emph{router\_fit} to produce the features the
router (Layer 2) trains on; for deployment, all nine experts are then refit
on the \emph{full} training split (disclosed simplification vs.\ full
$k$-fold cross-fitting -- the router itself is not refit at this stage,
trading a small amount of statistical efficiency for tractable runtime).

\subsection{Layer 2: graph mixture-of-experts router}
\label{sec:aegis_router}
A small gating network reads (a) each expert's probability, (b) a compact
graph/topology context (the two endpoint nodes' degree and the window's node
and edge counts, reused from the same windowed graphs as the RL environment's
observation), and (c) a handful of the flow's own raw feature values, and
outputs a top-$k$ ($k{=}3$) sparse softmax gate over the nine experts. The
router's predicted probability is the gate-weighted sum of expert
probabilities -- not a fresh classifier -- so routing only reweights the
pool's existing supervised signal per instance; it is trained (cross-entropy
against the true label, with a small gate-entropy regularizer discouraging
degenerate single-expert collapse) on the out-of-sample \emph{router\_fit}
scores described above, so it cannot learn to lean on an expert that has
merely overfit its own training data.

\subsection{Layer 3: tree-guided neural distillation}
\label{sec:aegis_distill}
A compact two-hidden-layer feed-forward student is trained against the
calibrated router-mixed ensemble's soft probability using four combined
losses: (i) temperature-scaled soft-label KL divergence against the teacher's
probability; (ii) a relational pairwise-similarity term that matches the
student's penultimate-layer embedding similarity structure, within a
minibatch, to the best single tree expert's leaf-index-embedding similarity
structure; (iii) a hard-class logit-margin term that pushes the minority
(attack) class's logit margin toward a target value; and (iv) class-imbalance
handling -- class-balanced focal loss in place of plain cross-entropy,
supervised contrastive learning on the student's embedding, and loss-weighted
resampling that oversamples the previous epoch's highest-loss examples. Each
of these four class-imbalance components is individually ablated
(Section~\ref{sec:results}) by disabling it alone, holding the other three
fixed.

\subsection{Layer 4: budgeted graph-RL investigation}
\label{sec:aegis_rl}
The POMDP action space is redesigned around a single structural change:
CLASSIFY\_BENIGN/CLASSIFY\_ATTACK are collapsed into one \textsc{Commit}
action, which always reads its label off the Layer-1/2 ensemble's calibrated
probability (Section~\ref{sec:aegis_calib}), thresholded at the policy's
currently selected operating point. The remaining nine actions are
\textsc{Abstain}, \textsc{RequestTelemetry}, \textsc{InspectNode},
\textsc{ExpandSubgraph}, \textsc{Escalate}, \textsc{Isolate},
\textsc{RaiseThreshold}, \textsc{LowerThreshold}, and \textsc{Defer}. This
guarantees the achievable classification quality is bounded below by the
calibrated ensemble's own quality regardless of how well or badly the RL
policy trains -- the exact failure mode identified in
Section~\ref{sec:method_history} is prevented structurally, not just
mitigated by a better reward. The observation adds, to the legacy
environment's masked-feature/graph-window state, the full per-expert
probability vector, the router's gate entropy and argmax-routed expert
identity, an approximate split-conformal prediction-set size, and the
calibrated probability itself as a class-risk score. The reward retains the
legacy terms (detection correctness, minority-class bonus, FN/FP penalty,
telemetry/inspection/subgraph/escalation cost, isolation cost/credit,
detection-delay bonus, budget-violation penalty) and adds a
calibration-violation penalty for committing immediately on a near-decision-boundary
probability without first acquiring information. Training uses the same
PPO-lite trainer (shared graph encoder, clipped surrogate objective, GAE)
as the legacy candidates, on the TRAIN split's environment, with greedy
evaluation on VAL and TEST. A third configuration, \textbf{Candidate C},
adds a fixed (non-learned) neuro-symbolic safety shield: if the current
flow's approximate conformal prediction set has size 2 (the calibrated
ensemble cannot exclude either class at $\alpha{=}0.05$) and the trained
policy chose \textsc{Commit}, the shield overrides the action to
\textsc{Escalate}. Candidates A (backbone only), B (+RL investigation), and C
(+safety shield) are compared using validation data only
(Section~\ref{sec:results}).

\textbf{Two properties of this environment, made explicit here rather than
left implicit.} First, \textsc{RequestTelemetry}/\textsc{InspectNode}/
\textsc{ExpandSubgraph} reveal previously-masked raw features and
graph-neighborhood statistics to the \emph{observation}, but the class-risk
score used by \textsc{Commit} is the fixed, precomputed ensemble probability
(Section~\ref{sec:aegis_calib}) -- it is never revised as telemetry is
acquired. The value of information from these three specific actions is
therefore analytically zero: they can only cost (telemetry/inspection/subgraph
fees, and a shrinking detection-delay bonus with each step taken before
committing), never help. \textsc{RaiseThreshold}/\textsc{LowerThreshold} are
not subject to this -- they change the operating point \textsc{Commit}
compares the (fixed) probability against, so they retain real effect. Second,
\textsc{Isolate}'s consequence (subsequent same-host flows auto-resolve to
the contained outcome) requires a real host identifier; CICIoT2023 and
UNSW-NB15 release none (Section~\ref{sec:datasets}), and the fallback used
elsewhere in this study for host-keyed bookkeeping degrades to a
per-row-unique index on both, meaning \textsc{Isolate} has \emph{zero}
containment effect on those two datasets -- only on ToN-IoT-Network, which
retains real source/destination IPs, does isolating a host actually contain
anything. We do not synthesize a proxy host key to manufacture sequential
structure that the released data does not support; isolation-as-containment
claims in Section~\ref{sec:results} are restricted to ToN-IoT-Network. An
earlier version of this environment's \textsc{Isolate} branch additionally
credited the containing host's isolation bonus \emph{on top of} the full
per-flow detection reward for the same flow (rather than in place of scoring
that flow only for bookkeeping) -- solving the resulting expected-reward
comparison analytically shows this made \textsc{Isolate} strictly dominate an
honest \textsc{Commit} for any flow with calibrated probability above
$p{=}2/3$, independent of whether isolating that specific flow was
informative. This was a reward-specification bug, not a deliberate design
choice, caught during review and fixed (the isolation credit is now scored
alone); Section~\ref{sec:results} reports both what this bug caused and what
changed once it was corrected, since the corrected behavior is what the
final numbers in this paper reflect.

\input{../tables/table_IIIb_aegis_pomdp.tex}

\subsection{Layer 5: calibration and low-FPR control}
\label{sec:aegis_calib}
The router-mixed ensemble probability is calibrated on the VALIDATION split
only using four methods evaluated side by side: temperature scaling
(single-parameter logistic recalibration), isotonic regression, Platt
scaling (logistic regression on the logit-transformed score), and
class-wise/FPR-budget threshold selection; a split-conformal procedure
additionally computes a threshold controlling benign FPR at a target
$\alpha$ on held-out validation benign flows. TEST data is scored under
whichever calibration/threshold was selected on VAL -- never refit or
threshold-tuned on TEST. All four calibration variants are compared in
Section~\ref{sec:results} by expected calibration error (ECE), Brier score,
negative log-likelihood, and resulting FPR/FNR/alert volume.

\subsection{Algorithm and complexity}
Algorithm~\ref{alg:aegis} summarizes training. Inference cost at deployment
is dominated by whichever configuration is selected: the full nine-expert
pool plus router (highest accuracy, highest cost), the distilled student
(lowest cost, reduced accuracy, Section~\ref{sec:results}), or the RL policy
layered on top of either (adds one small actor-critic forward pass per
investigation step, bounded by the per-episode telemetry budget). Parameter
counts, latency, throughput, and memory for every configuration are measured
directly rather than estimated (Section~\ref{sec:results}).

\begin{algorithm}[t]
\caption{AegisGraph-X training (one dataset)}
\label{alg:aegis}
\begin{algorithmic}[1]
\State Leakage-safe split into TRAIN/VAL/TEST (Section~\ref{sec:datasets})
\State Split TRAIN $\to$ expert\_fit (70\%), router\_fit (30\%), contiguous
\State Fit 9 experts on expert\_fit; score out-of-sample on router\_fit
\State Train MoE router on router\_fit's out-of-sample scores + graph context
\State Refit all 9 experts on full TRAIN (router held fixed)
\State Score VAL/TEST through refit experts $\to$ router $\to$ ensemble probability
\State Fit calibration bundle (temperature/isotonic/Platt/conformal) on VAL
\State \textbf{Candidate A} $\gets$ calibrated ensemble (Layers 1/2/5 only)
\State Train tree-guided distilled student on TRAIN (Layer 3)
\State Build AegisPOMDPEnv (Layer 4) from calibrated ensemble + expert probs
\State Train PPO investigation policy on TRAIN env $\to$ \textbf{Candidate B}
\State \textbf{Candidate C} $\gets$ Candidate B + conformal safety shield
\State Select final method among A/B/C using VAL-only criteria (Section~\ref{sec:results})
\end{algorithmic}
\end{algorithm}

\section{Experimental Setup}
\label{sec:setup}

\textbf{Datasets.} CICIoT2023, ToN-IoT-Network, UNSW-NB15 (primary),
RT-IoT2022 (supplementary transfer-only), as justified in
Section~\ref{sec:datasets}.

\textbf{Baselines.} The same 37 classical/deep/graph/RL baselines validated in
our prior study are reused unchanged (12 classical/tree, 9 deep, 5 graph, 11
RL/decision -- full list in Section~\ref{sec:method_history}'s predecessor
manuscript and \texttt{code/baselines/}), plus our own prior offline-safe-RL
design kept as baseline \#39, plus three AegisGraph-X ablation variants
reported alongside the full method as baselines \#40-42 (best single tree
expert alone, uncalibrated MoE ensemble, distilled student) -- 41 comparison
points in total against AegisGraph-X (\#43). \emph{Scope note:} the task's
suggested expansion to a 43-baseline suite (adding an RBF-kernel SVM,
TabTransformer, and SAINT as new baselines) was not implemented in this
revision, to concentrate the available time budget on the new method itself
rather than re-running an already-large, already-validated baseline suite;
this is a disclosed tractability decision, not a hidden gap, and the
already-implemented 37 span the same classical/deep/graph/RL families those
three additions would have.

\textbf{Hyperparameter tuning.} Every AegisGraph-X component's thresholds are
selected on the VALIDATION split only (Section~\ref{sec:aegis_calib}); the
RL investigation layer's reward weights are evaluated under two explicit
modes -- common fixed weights, and a small validation-only grid search per
dataset (Section~\ref{sec:results}) -- never tuned or selected using TEST
data. The legacy Candidate-3 baseline retains its own previously validated,
validation-only Optuna-tuned configuration unchanged.

\textbf{Metrics.} Accuracy, balanced accuracy, macro/weighted F1, attack and
minority recall, per-class precision/recall/F1, ROC-AUC, PR-AUC, FPR, FNR,
cumulative reward, detection delay, telemetry cost, escalation/abstention/isolation
rate, budget-violation rate, latency, throughput, parameter count, and memory.

\textbf{Statistical testing.} Repeated seeds (full pairwise table in the
code repository's \texttt{results/csv/}), paired $t$-test and Wilcoxon
signed-rank test with Holm correction, and effect
size, computed only over metrics logged from real runs.

\textbf{Hardware.} NVIDIA RTX 5090 (24\,GB), 24-core CPU, PyTorch with CUDA
12.8, PyTorch Geometric 2.7.

\section{Results and Discussion}
\label{sec:results}

\input{../tables/table_IV_candidate_screening.tex}

\subsection{Historical context: the legacy candidate screening}
Table~\ref{tab:candidate_screening} is retained unchanged from our prior
study: it screened ten RL-as-classifier candidates and selected Candidate 3
(Offline Safe Graph RL, $S{=}0.538$) purely on logged validation evidence.
Section~\ref{sec:method_history} summarizes why this design was not retained
as AegisGraph-X's headline method; Candidate 3 is kept as baseline \#39 in
every comparison below.

\subsection{Validation-only selection among Candidates A/B/C}
\label{sec:results_selection}
\input{../tables/table_aegis_selection.tex}
Table~\ref{tab:aegis_selection} reports the outcome of the Part-3 selection
protocol: each candidate's gated score is its mean validation macro-F1
multiplied by a continuous gate on its \emph{worst-dataset} minority recall
(floor $0.30$) -- the same gating philosophy used to fix the legacy
Candidate-3 tuning collapse (Section~\ref{sec:method_history}), applied here
to candidate \emph{selection} rather than hyperparameter tuning.
\textbf{Candidate A (backbone only, no RL) is selected}: on CICIoT2023's
extreme attack-majority class balance, the trained RL policy (Candidates B
and C) collapsed to an always-\textsc{Isolate} policy under the common
default reward weights (minority/benign recall $0$, FPR $1.0$) -- a genuine,
disclosed negative result, not a hidden one, and exactly the kind of collapse
the selection protocol's minority-recall gate is designed to catch. On the
other two datasets, B and C do not collapse (UNSW-NB15 minority recall
$0.92$--$0.93$; ToN-IoT-Network $0.99$), but neither improves classification
over A once collapse is accounted for. \textbf{We therefore report explicitly:
the RL investigation layer, at the common reward weights and training budget
evaluated here, does not improve classification accuracy over the backbone
alone, and on one of three datasets it substantially degrades it} -- this is
the honest outcome required by the task's wording policy, not a result we
searched for. Section~\ref{sec:results_rewardmodes} reports whether a
validation-only reward-weight search (Mode 2) recovers B/C's classification
quality without touching test data.

\input{../tables/table_V_aegis_binary.tex}

\subsection{Binary detection performance}
Table~\ref{tab:aegis_binary} reports AegisGraph-X's (Candidate A, the
selected method) TEST-split binary detection performance. Over ten seeds at a pinned training budget, macro-F1 reaches
$0.940\pm0.002$ (UNSW-NB15), $0.997\pm0.002$ (ToN-IoT-Network), and
$0.949\pm0.007$ (CICIoT2023) -- a categorical improvement over the legacy
method's $0.465$--$0.679$, and no longer 20--53 points behind the strongest
tree baseline on any dataset. The residual gap to that baseline is
$1.2$ points on CICIoT2023 and under $0.1$ points on the other two
(Section~\ref{sec:results_equivalence}).

\subsection{Comparison with the strongest tree/boosting baseline}
\label{sec:results_equivalence}

AegisGraph-X is compared against the single strongest tree/boosting baseline
per dataset (XGBoost, LightGBM, and GradientBoosting respectively). An
earlier version of this manuscript described the outcome as a statistical
\emph{tie} on the grounds that a Holm-corrected paired $t$-test failed to
reject the null, and reported it in a table structured around that framing.
That inference was wrong in kind, not merely in degree: failure to reject is
not evidence of equivalence, and at three seeds the test has almost no power
to detect differences of the size at issue. We therefore replace that table
and its framing entirely with a two one-sided tests (TOST)
procedure~\cite{schuirmann1987tost,lakens2017equivalence} against a
pre-specified margin, which asks the question the tie claim was informally
reaching for: is the difference small enough to be operationally irrelevant?

\input{../tables/table_tost_equivalence.tex}

Table~\ref{tab:aegis_tost} reports the result over ten seeds at a single
pinned row budget (Section~\ref{sec:limitations}). The outcome separates the
three datasets cleanly rather than uniformly. On UNSW-NB15 and
ToN-IoT-Network, \textbf{equivalence holds at the pre-specified
$\delta = 0.005$ with substantial margin to spare}: the smallest attaining
margins are $\delta^{*} = 0.0010$ and $0.0011$ respectively, so the two
methods are indistinguishable to within roughly one tenth of a macro-F1
point. On UNSW-NB15 the paired difference is in fact nominally positive
($+0.0004$).

\textbf{CICIoT2023 does not clear the margin, and this is a real deficit
rather than a power artifact.} The paired difference is $-0.0117$, the
$90\%$ confidence interval $[-0.0148, -0.0087]$ excludes zero entirely, and
the TOST $p$-value at $\delta = 0.005$ is $0.999$. Equivalence is attained
only at $\delta = 0.02$, with $\delta^{*} = 0.0148$. AegisGraph-X is
reliably worse than gradient boosting on CICIoT2023 by approximately $1.2$
macro-F1 points.

The defensible claim is therefore \textbf{equivalence within a tenth of a
point on two corpora and a measurable deficit of $1.2$ points on the third},
which is both more precise and less flattering than the tie language it
replaces. We note that the earlier three-seed analysis was wrong in both
directions at once: it understated the CICIoT2023 deficit (reporting
$-0.0066$ against the true $-0.0117$) while also understating the UNSW-NB15
and ToN-IoT-Network agreement (reporting equivalence margins of $0.82$ and
$0.28$ points against the true $0.10$ and $0.11$). Both errors trace to the
same two causes -- three seeds, and a training-set size that differed
between the method and its baselines -- and both are corrected here.

\input{../tables/table_XVII_aegis_win_tie_loss.tex}
\input{../tables/table_XVI_aegis_stat_sig.tex}

\begin{figure}[t]
\centering
\includegraphics[width=\linewidth]{../figures/fig_aegis_win_tie_loss.png}
\caption{Statistically superior / tied / inferior counts vs.\ 41 other models, per dataset.}
\label{fig:aegis_win_tie_loss}
\end{figure}

\subsection{Full benchmark comparison and win/tie/loss}
Table~\ref{tab:aegis_win_tie_loss}/\ref{tab:aegis_stat_sig} and
Figure~\ref{fig:aegis_win_tie_loss} report the
complete comparison against all 41 other rows in the benchmark suite (37
original baselines, the legacy RL method \#39, and proposed-variant ablations
\#40-42), Holm-corrected paired $t$-test over the 3 seeds common to
AegisGraph-X and every other model. \textbf{AegisGraph-X is statistically
superior to 11 baselines on UNSW-NB15, 4 on ToN-IoT-Network, and 13 on
CICIoT2023 (28 statistically-superior comparisons in total across the three
datasets), is statistically tied with the remaining classification-competitive
baselines, and is statistically inferior to exactly one comparison across all
three datasets: the OracleUpperBound baseline on UNSW-NB15} (which reads the
ground-truth label directly and is never presented as a deployable method).
\textbf{AegisGraph-X is never significantly beaten by any real (non-oracle)
baseline on any dataset} -- a qualitative change from the legacy method,
which lost 19 significant comparisons to classical/tree/deep/graph baselines
and won none significantly.

This last statement requires a caveat that we did not give it in the earlier
version. ``Never significantly beaten'' is a claim about non-rejection, and
with three seeds and Holm correction~\cite{holm1979} across 41 comparisons
per dataset it is close to guaranteed a priori: clearing the most stringent
Holm threshold requires $t \gtrsim 25$ at two degrees of freedom, so the
procedure is far better at failing to detect losses than at detecting them.
The result is therefore weak evidence of parity and should not be read as a
property of the method. The equivalence analysis in
Section~\ref{sec:results_equivalence} is the load-bearing comparison against
strong baselines; the win counts below are informative mainly against the
weak baselines where the gaps are large. The statistically-superior comparisons span every
non-tree family: weaker classical baselines (logistic regression, GaussianNB,
LinearSVM), weaker deep/graph baselines (FT-Transformer, GCN, HypergraphNN),
and -- notably -- \emph{every} RL/decision-policy baseline evaluated on the
datasets where the comparison reaches significance, including PPO, A2C,
offline CQL, static-threshold, and both greedy-acquisition baselines, and,
on CICIoT2023, the legacy Offline Safe Graph RL method itself
($p_{\text{Holm}}{=}0.009$).

We also computed the paired Wilcoxon signed-rank test and Cohen's $d$ for
every comparison, as Section~\ref{sec:setup} states we would, and report
here that the Wilcoxon result does not corroborate the paired $t$-test at
this seed count: none of the 28 comparisons the $t$-test calls
statistically superior remains significant under Wilcoxon with Holm
correction ($n=3$ paired seeds gives the signed-rank statistic too few
attainable values to reach a two-sided $p<0.05$ regardless of effect size).
Cohen's $d$ for those same 28 comparisons ranges from $18$ to $318$, which
is not a meaningful effect-size scale at $n=3$ -- it reflects near-zero
pooled variance across three seeds as much as it reflects a real
difference. We report both numbers rather than the more flattering one:
the $t$-test superiority claims in this section should be read as
consistent with a large nominal gap, not as independently corroborated by
a second test, and the seed-count limitation stated at the start of this
section applies here with full force.

\subsection{Ablation}
\label{sec:results_ablation}
\input{../tables/table_XI_aegis_ablation.tex}

\begin{figure}[t]
\centering
\includegraphics[width=\linewidth]{../figures/fig_aegis_ablation.png}
\caption{Component ablation, TEST macro-F1 (mean over 3 datasets $\times$ 10 seeds).}
\label{fig:aegis_ablation}
\end{figure}

Table~\ref{tab:aegis_ablation} reports the full component ablation (TEST
split, mean over 3 datasets $\times$ 10 seeds, $n = 30$ paired
observations; Figure~\ref{fig:aegis_ablation} plots the same data). At three
seeds the two headline effects below were nominal; at ten, one of them is
significant and the other is a clean null, and both run against the paper's
own narrative expectations rather than for it.

\textbf{(1) The nine-expert pool is significantly worse than its own best
member.} Using the single best expert alone reaches macro-F1 $0.9652$
against the full pool's $0.9621$ -- a paired difference of $+0.0031$ in
favour of the ablation, $t = 2.87$, $p = 0.008$ over $n = 30$. This is not
the weaker statement that the ensemble is redundant; the ensemble actively
costs accuracy relative to picking one of its constituents. We report it as
the clearest evidence against Layer 1 as specified.

\textbf{(2) The learned router is indistinguishable from an unweighted
average.} Equal-weight averaging of all nine experts gives $0.9617$ against
the router's $0.9621$, a paired difference of $-0.0004$ with $p = 0.90$.
This is a genuine null rather than an underpowered one: at $n = 30$ the
interval excludes any effect large enough to matter operationally. We
therefore do not claim the learned router improves classification accuracy
over a much simpler ensembling rule at this training scale. Its value, if
any, lies in the per-instance routing signal it exposes to the RL layer
(argmax expert identity, gate entropy), not in raw ensemble accuracy. The
result is consistent with the stacked-generalization
literature~\cite{wolpert1992stacked}, where learned combiners routinely fail
to beat simple averaging at modest sample sizes.

Dropping the graph branch ($0.9624$) or the temporal branch ($0.9615$) costs
less than a tenth of a macro-F1 point relative to the full pool, indicating
the six tree/boosting experts carry essentially all of the classification
signal on these three datasets -- consistent with the long-standing finding
that gradient-boosted trees are hard to beat on tabular
data~\cite{grinsztajn2022trees,shwartzziv2022tabular}.

\textbf{The tree-guided distillation ablation is a second negative result,
and the larger seed count makes it worse than we previously reported.} The
full compact student (focal loss + supervised contrastive learning +
hard-negative mining all enabled) reaches macro-F1 $0.8417$, a substantial
and expected drop from the nine-expert pool ($0.9621$) given its
$\sim$8K-parameter budget. Only focal loss earns its place, and even that
marginally: removing it costs $0.0305$ ($p = 0.054$).
\textbf{Removing hard-negative mining \emph{raises} macro-F1 by $0.0449$
($p = 0.008$), and the plain-KD-only student with all three techniques
disabled reaches $0.8978$ -- $5.6$ points above the full student
($p = 0.001$).} Removing contrastive learning also raises the mean, by
$0.0238$, though at $p = 0.158$ we do not claim that individual effect is
resolved. The bundle as a whole is significantly harmful, so the honest
summary is that Layer 3 would be better specified as plain knowledge
distillation~\cite{hinton2015distilling} with focal
loss~\cite{lin2017focal} alone. We retain all three as ablatable,
documented components rather than silently dropping them, but we do not
claim contrastive learning or hard-negative mining is validated here.

One component does earn its place. Disabling calibration and thresholding
the raw ensemble probability at $0.5$ costs $0.0028$ macro-F1
($p = 0.040$) -- a small but statistically resolved benefit, and the only
one of AegisGraph-X's additions that survives its own ablation at ten seeds.
We note it explicitly because the rest of this section is negative, and the
distinction between calibration (which helps) and the conformal FPR-control
variant built on top of it (which, per
Section~\ref{sec:results_lowlabel_robust}, does not measurably change alert
volume) is easy to lose.

\subsection{Dual-mode RL evaluation (common vs.\ validation-tuned reward)}
\label{sec:results_rewardmodes}
\input{../tables/table_IX_aegis_common_reward.tex}
\input{../tables/table_X_aegis_tuned_reward.tex}
Table~\ref{tab:aegis_common_reward} (Mode 1: the same fixed default reward
weights used throughout Section~\ref{sec:results_selection}) and
Table~\ref{tab:aegis_tuned_reward} (Mode 2: a small validation-only
per-dataset reward-weight grid search over the isolation-cost/credit and
false-positive terms, Section~\ref{sec:aegis_rl}) are reported separately, and
raw cumulative reward is never compared across the two modes or across
datasets, since the weights differ (Part 7 of the task protocol). This
dual-mode comparison is a smaller, exploratory pass (3 seeds, the
$50$K-row budget used before the training-budget-parity correction in
Section~\ref{sec:limitations}) than the ten-seed, pinned-row-budget
evidence that Section~\ref{sec:results_selection} actually selects on, and
the two disagree in an instructive way: at this smaller scale, only $1$ of
the $18$ (dataset, seed, mode) training runs collapses to the degenerate
always-\textsc{Isolate} policy (minority recall $<0.05$; ToN-IoT-Network,
seed 1, Mode 2), telemetry spend and isolation rate are both $0$ in every
other run, and CICIoT2023 itself never collapses at this scale.
\textbf{The primary evidence for Candidate A's selection is not this
dual-mode pass but the ten-seed, pinned-row-budget result in
Section~\ref{sec:results_selection}}, where the same always-\textsc{Isolate}
collapse (minority recall $<0.05$, isolation rate $\approx0.90$) recurs in
$9$ of $10$ seeds on CICIoT2023 under the common default reward weights,
while UNSW-NB15 and ToN-IoT-Network never collapse and instead settle into a
substantial, non-degenerate isolation rate ($0.58$ and $0.54$ respectively)
that does not cost minority recall. Reading the two studies together, the
honest characterization is scale-dependence, not the bimodal
seed-to-seed instability an earlier version of this section claimed: the
same reward specification and the same fixed ISOLATE-credit code
(Section~\ref{sec:aegis_rl}) produce a training regime that is far more
prone to the CICIoT2023 collapse at the larger, production-realistic row
budget than at the smaller one, for reasons this study does not isolate
(more of the same class-imbalance-heavy data to fit an on-policy
advantage estimate against, more updates per fixed episode-count schedule,
or simply a different sampling of an already-unstable optimization
landscape). We do not have a controlled sweep across row budgets to
distinguish these, and flag it as the most direct follow-up this finding
implies. Where the RL layer does not collapse, its decision-cost behavior
(telemetry cost, escalation/isolation rate, detection delay --
reward-scale-independent columns in
Tables~\ref{tab:aegis_common_reward}--\ref{tab:aegis_tuned_reward}) is
reported for completeness, but we do not build a deployment recommendation
on a policy that fails outright in $9$ of $10$ seeds on one of three
evaluation datasets at the row budget the rest of this paper uses.

\subsection{The analytic decision-theoretic baseline}
\label{sec:results_analytic}
\input{../tables/table_aegis_analytic_baseline.tex}
Every earlier version of this paper flagged the same gap: we compare the RL
layer against \emph{no} RL, but never against the closed-form,
expected-cost-minimizing policy the same calibrated posterior and the same
reward table admit analytically, with no training at all
(Section~\ref{sec:aegis_rl}). We compute that policy directly -- for each
flow, the argmax over the five terminal actions' expected reward under
$p=$~AegisGraph-X's own calibrated ensemble probability -- and evaluate it
on the identical ten-seed, pinned-row-budget test data used throughout
Section~\ref{sec:results_selection}. Table~\ref{tab:aegis_analytic} reports
the result, paired by seed against Candidate A.

\textbf{The analytic policy never isolates on any of the three datasets
(isolation rate $0.000$ everywhere, including CICIoT2023), and it
significantly outperforms Candidate A on two of three datasets} ($+0.0222$
macro-F1 on UNSW-NB15 and $+0.0219$ on CICIoT2023, both $p<0.001$; $+0.0006$
on ToN-IoT-Network, $p=0.112$, not significant). This resolves the open
question Section~\ref{sec:limitations} raised and could not previously
settle. The isolate-reward fix (Section~\ref{sec:aegis_rl}) removed the
$p=2/3$ dominance crossover that made \textsc{Isolate} strictly better than
an honest \textsc{Commit} above that confidence threshold; under the
corrected reward table, the closed-form optimum simply never isolates,
anywhere, on any of the three datasets -- \textbf{the CICIoT2023 collapse
the trained RL policy exhibits in $9$ of $10$ seeds is therefore not a
property of the reward specification the policy was trained against. The
analytic optimum, computed from that identical specification, does not
collapse.} The most economical remaining explanation is that PPO training
itself finds and gets stuck in a spurious always-\textsc{Isolate} local
optimum that the true expected-reward landscape does not actually contain,
rather than that the landscape rewards collapse and the policy is simply
finding it. We had flagged, in earlier drafts of this section, that our
experiments could not distinguish a reward-specification explanation from a
training explanation for the CICIoT2023 collapse; this comparison
distinguishes them, and the training explanation is the one the evidence
supports.

The analytic policy's advantage over Candidate A on UNSW-NB15 and CICIoT2023
has a second, independent explanation, and both likely contribute: Candidate
A's \textsc{Commit} threshold is drawn from a five-point preset grid
($\{0.3, 0.4, 0.5, 0.6, 0.7\}$, Section~\ref{sec:aegis_rl}), while the
closed-form policy's implicit commit/no-commit boundary is continuous and
falls at approximately $p \approx 0.26$ given the reward weights
(Section~\ref{sec:aegis_rl}) -- below the grid's minimum. Candidate A cannot
reach this operating point regardless of training; the analytic policy can
by construction. We do not separate the two effects (reward-table-optimal
action choice vs.\ threshold granularity) in this comparison and flag doing
so as a natural follow-up, but note that fixing the threshold-granularity
issue alone -- widening or removing the preset grid -- is a small, low-risk
change relative to redesigning the RL layer, and is the more promising
near-term fix implied by this result.

We report this comparison as a genuine addition to the paper's evidence
base, not a validation of AegisGraph-X's RL layer: the finding is that a
policy which does not learn anything, is not sequential, and never spends a
telemetry budget outperforms both the trained RL policy and, on two of
three datasets, the backbone's own default operating point. If anything,
this sharpens rather than softens the paper's central negative result about
Layer 4.

\subsection{Low-label, robustness, and calibration}
\label{sec:results_lowlabel_robust}
\input{../tables/table_XII_aegis_low_label.tex}
\input{../tables/table_XIII_aegis_robustness.tex}
\input{../tables/table_XIV_aegis_calibration.tex}
\textbf{Low-label} (Table~\ref{tab:aegis_low_label}, seed 0, full backbone
refit at each fraction): AegisGraph-X degrades gracefully rather than
falling off a cliff as training labels shrink from 100\% to 1\%: macro-F1
$0.933\rightarrow0.891$ (UNSW-NB15), $0.997\rightarrow0.977$
(ToN-IoT-Network), and $0.950\rightarrow0.783$ (CICIoT2023) --
CICIoT2023 is the most label-hungry of the three, consistent with it having
the largest and most heterogeneous attack-class composition of the three
primary datasets even under the binary framing used here.

\textbf{Robustness} (Table~\ref{tab:aegis_robustness}, 2 seeds): no
perturbation leaves AegisGraph-X unaffected, and we report the most fragile
result rather than lead with the most flattering one -- \textbf{ToN-IoT-Network,
despite the highest clean-condition macro-F1 of the three datasets ($0.996$),
is also the \emph{most} fragile to feature-noise perturbation, dropping to
macro-F1 $0.38$--$0.44$ under moderate/high additive noise}, a larger
relative collapse than either UNSW-NB15 ($0.94\rightarrow0.67$--$0.76$) or
CICIoT2023 ($0.95\rightarrow0.56$--$0.65$) show under the same perturbation.
Missing-feature injection is damaging but less catastrophic (macro-F1 drops
of $0.09$--$0.36$ across datasets at 50\% missingness); class-imbalance
stress and 5\%-label-noise are the mildest perturbations tested (drops
$\le\!0.07$ on all three datasets). We do not claim AegisGraph-X is robust
to feature corruption -- these results show it is not, and that its
near-ceiling clean-condition performance on ToN-IoT-Network specifically
does not carry over to noisy conditions.

\textbf{Calibration} (Table~\ref{tab:aegis_calibration}, 1 seed, mean over 3
datasets): the router-mixed ensemble's raw probability is already
reasonably well-calibrated before any post-hoc correction (uncalibrated ECE
$0.0038$, Brier $0.0148$) -- expected, since tree/boosting ensembles tend to
produce reasonably calibrated probabilities on their own. Isotonic
regression gives the largest additional improvement (ECE $0.0023$), Platt
scaling and temperature scaling are intermediate, and the split-conformal
FPR-control variant (fit to target $\alpha{=}0.05$ on validation benign
flows) achieves comparable FNR/alert-volume to the other variants while
providing a distribution-free FPR guarantee. All differences among the four
calibrated variants are small relative to the uncalibrated baseline, so we
do not claim calibration is a major lever for this method on these three
datasets, only that it is available and validated.

\input{../tables/table_VII_multiclass.tex}

Multiclass evaluation (Table~\ref{tab:multiclass}) is reused unchanged from
our prior study as a classical-baseline reference point; AegisGraph-X's
router and RL layers remain binary-only in this revision, the same
disclosed limitation as before (Section~\ref{sec:limitations}), since
extending the MoE router and calibration layers to a native multiclass
target is a substantial additional engineering effort we did not complete
within this revision's time budget.

\input{../tables/table_XV_aegis_runtime.tex}

\subsection{Runtime and deployment}
Table~\ref{tab:aegis_runtime} measures inference latency, training time, and
parameter count directly (never estimated) for four deployment
configurations, over 10 seeds. The single strongest tree expert alone is
fastest among the full-accuracy options ($16.3\mu$s/sample); the full
nine-expert pool plus MoE router costs $\sim\!9\times$ more per sample
($142\mu$s) for the accuracy reported above; the RL investigation layer's row
records the same latency as the ensemble row to three significant figures,
which we read as the policy forward pass being too small to resolve against
the graph-reconstruction cost rather than as a measurement of zero -- we do
not claim the RL layer is free at inference time. Absolute latencies here are
wall-clock measurements on a shared, long-running research machine and
fluctuate across runs by roughly an order of magnitude in our own repeated
measurements; the relative ordering (student $\ll$ single tree $<$ full
ensemble) is stable across every run we have taken, but the specific
multipliers should be read as indicative, not precise. \textbf{The distilled
student is $\sim\!200\times$ faster than the single tree expert and roughly
three orders of magnitude faster than the full ensemble} ($0.076\mu$s/sample)
at a real, disclosed accuracy cost (macro-F1 $0.842$ vs.\ $0.962$, 10-seed
ablation). This is a measured latency/accuracy tradeoff for
resource-constrained settings, not evidence that the compact student matches
the full method's accuracy.

One clarification about model size, since an earlier version of this text
invited a misreading. The student has $\sim\!8.1$K parameters and the RL
policy $\sim\!9.9$K, so the student is $18\%$ smaller \emph{than the policy
network} -- but the tree ensembles have no parameter count in the same sense,
and Table~\ref{tab:aegis_runtime} accordingly leaves that field empty for
them. The $740\times$ figure is therefore a latency result, not a compression
result, and the parameter counts should not be read as quantifying
compression against the expert pool.

\emph{Full per-model/per-dataset detail for every result summarized above is
provided as generated LaTeX/CSV in the code repository's \texttt{tables/} and
\texttt{results/csv/} directories rather than inlined here, to respect the
page limit.}

\section{Limitations}
\label{sec:limitations}

Dataset selection followed explicit, predeclared criteria but does not imply
universal dominance across all IDS corpora; three of the twelve inspected
dataset families are used as primary evidence, and cross-dataset transfer to
RT-IoT2022 was not re-evaluated for AegisGraph-X in this revision (a
disclosed scope reduction). The POMDP environment is reconstructed from
static, already-labeled benchmark data rather than a live SOC feed:
escalation is modeled as reliably resolving to the ground-truth label at a
fixed cost, and host isolation's effect on subsequent flows is simulated
within an episode rather than observed from a real network response, so
simulated response actions should not be equated with live deployment
behavior. \textbf{The MoE router's classification accuracy is not
statistically distinguishable from a much simpler unweighted expert average
or the single best tree expert alone} (Section~\ref{sec:results}) -- we do
not claim the learned router is validated as a classification-accuracy
improvement, only as an infrastructure that exposes useful per-instance
signal (routed-expert identity, gate entropy) to the RL layer. \textbf{The
budgeted RL investigation layer, at the common default reward weights,
collapsed to a degenerate always-isolate policy on CICIoT2023} and did not
improve classification accuracy over the backbone alone on the other two
datasets either; validation-only reward-weight tuning (Mode 2,
Section~\ref{sec:results_rewardmodes}) changes this only partially. We do not
claim the RL investigation layer is validated as a classification-accuracy
improvement in this study; its structural role (Section~\ref{sec:aegis_rl})
is to bound classification quality below by the backbone's own quality while
adding sequential decision-cost control, and its decision-utility value,
where present, is reported separately from classification accuracy
throughout (Part 16 wording policy). \textbf{The tree-guided distillation's
hard-negative-mining component did not improve the compact student on
average across the three datasets} (Section~\ref{sec:results}), an honest
negative result for one of Layer 3's four components. The baseline suite
was not expanded to the task's suggested 43-model list (RBF-SVM,
TabTransformer, SAINT) in this revision, a disclosed tractability decision
(Section~\ref{sec:setup}). Multiclass evaluation remains a classical-baseline
reference point only, as in our prior study; AegisGraph-X's router,
calibration, and RL layers were not extended to a native multiclass target.
Physical deployment latency/memory are measured on the reported hardware
only and have not been validated on resource-constrained edge devices.
\textbf{Seed count remains uneven across experiments.} The primary
comparisons in Sections~\ref{sec:results_equivalence}
and~\ref{sec:results_ablation} use ten seeds, which is sufficient to resolve
both the equivalence bounds and the two significant ablation effects. The
supporting experiments are weaker and were not re-run: the low-label study
(Table~\ref{tab:aegis_low_label}) uses a single seed, the calibration
comparison (Table~\ref{tab:aegis_calibration}) uses a single seed while
reporting differences in the fourth decimal, and the robustness study
(Table~\ref{tab:aegis_robustness}) uses two. Comparative claims resting on
those three tables should be read as indicative only. We also restrict the
ten-seed budget to the classical/tree family and the proposed method, since
those support the equivalence and ablation \emph{nulls} where power is the
argument; the deep, graph, and RL baseline families retain three seeds,
where the margins are large and the conclusions are superiority claims
rather than null results.

\textbf{Training-budget parity between the method and its baselines was not
verified in the original run, and correcting it changed the results.} The
experiment drivers used different row budgets -- the AegisGraph-X driver
applied a fixed cap while the baseline driver used per-dataset defaults --
so the head-to-head comparison in Table~\ref{tab:aegis_tost} was, as
originally executed, confounded with training-set size. The split manifests
reproduced with the submitted results also did not match the row counts
reported in Table~\ref{tab:dataset_stats}. All results reported here are
regenerated under a single pinned row budget applied identically to every
model, with split manifests emitted from the same run that produced the
tables. The correction moved the comparison in both directions: the
UNSW-NB15 and ToN-IoT-Network equivalence bounds tightened by roughly a
factor of eight, while the CICIoT2023 deficit nearly doubled. We record this
because the original numbers were not merely imprecise -- they were
uninterpretable, and their errors did not share a sign.

\textbf{The budgeted acquisition layer is unexercised, not merely
unhelpful.} Tables~\ref{tab:aegis_common_reward} and~\ref{tab:aegis_tuned_reward}
(the three-seed dual-mode study) show telemetry cost and isolation rate both
exactly zero on every dataset under both reward modes -- the policy never
spends its acquisition budget and never isolates a host at this scale. This
is consistent with Section~\ref{sec:aegis_rl}'s analytic argument that
\textsc{RequestTelemetry}/\textsc{InspectNode}/\textsc{ExpandSubgraph} have
zero value of information by construction (\textsc{Commit} reads a fixed,
precomputed probability those actions never revise), so a
reward-maximizing policy correctly learns not to pay for them. Isolation is
a different case: it does carry real expected value once the fixed-reward
bug is corrected (Section~\ref{sec:aegis_rl}), and at the larger, ten-seed,
pinned-row-budget scale used for the primary comparisons
(Section~\ref{sec:results_selection}) the policy does use it substantially
-- isolation rate $0.58$ (UNSW-NB15) and $0.54$ (ToN-IoT-Network), without
cost to minority recall on either dataset -- and, on CICIoT2023, uses it
almost exclusively ($0.90$) in a way that does cost minority recall in $9$
of $10$ seeds. The honest reading is therefore two separate findings, not
one: the three acquisition actions are unexercised everywhere because the
environment gives them no information to act on, which is a property of the
design, not of training; the isolation action is exercised, sometimes to
good effect and sometimes catastrophically, in a manner that depends on
dataset and training scale in ways this study does not fully explain.

\textbf{The analytic decision-theoretic baseline (Section~\ref{sec:results_analytic})
resolves part of this paper's central open question, and sharpens rather
than softens the negative result.} Once Layer 5 supplies a calibrated
posterior and the reward specifies an explicit cost structure, the
cost-minimizing action is computable in closed form by expected-cost
minimization~\cite{howard1966voi,wald1945sequential}, with no training. That
policy never isolates on any of the three datasets and significantly
outperforms Candidate A on two of three -- so the CICIoT2023 collapse the
trained RL policy exhibits in $9$ of $10$ seeds is attributable to PPO
training failing to find a landscape optimum that does not itself reward
collapse, not to the reward specification rewarding collapse. This narrows,
but does not close, the earlier uncertainty: we still do not know why PPO
finds this spurious optimum at the larger row budget and mostly does not at
the smaller one (Section~\ref{sec:results_rewardmodes}), only that the
optimum it finds is not the one the reward table actually defines.

\textbf{The structural guarantee is a design property, not a result.} That
classification quality cannot fall below the calibrated ensemble's is true by
construction, since COMMIT reads the ensemble's label
(Section~\ref{sec:aegis_rl}). It should not be read as a finding, and we note
that it bounds only the metric we report while leaving unbounded the
decision-cost behaviour that Tables~\ref{tab:aegis_common_reward}
and~\ref{tab:aegis_tuned_reward} show degrading.

\textbf{Robustness undercuts the deployment framing.} Relative to its own
two-seed clean baseline (Table~\ref{tab:aegis_robustness}; note that this
study was not re-run at ten seeds or at the pinned row budget, so its clean
figures differ slightly from Table~\ref{tab:aegis_binary}), moderate feature
noise drives ToN-IoT-Network macro-F1 from $0.996$ to $0.38$--$0.44$,
UNSW-NB15 from $0.938$ to $0.67$--$0.76$, and CICIoT2023 from $0.952$ to
$0.56$--$0.65$. In intrusion detection, feature perturbation is the threat
model rather than an incidental robustness question, so these numbers bound
the operational reading of the clean-condition results more tightly than the
clean results themselves suggest.

\textbf{Low-FPR control is not demonstrated.} The conformal FPR-control
variant reduces alert volume from $7948.1$ to $7934.9$ per $10{,}000$ flows
at $\alpha = 0.05$ (Table~\ref{tab:aegis_calibration}), a change of $0.2\%$.
On these corpora, whose attack prevalence ranges from $64\%$ to $98\%$, the
alert-volume metric has little headroom and the ``low-FPR control'' framing
is not supported by the measurement. Establishing it would require
evaluation at a realistic base rate, which these datasets do not provide.

\textbf{Deployment and harm.} This paragraph concerns Candidates B and C
(the RL investigation layer), not the selected Candidate A, which never
isolates a host at all. At the ten-seed, pinned-row-budget scale, the
evaluated RL policies isolate $54$--$90\%$ of hosts depending on dataset
(Section~\ref{sec:results_rewardmodes}); at the smaller three-seed dual-mode
scale the same policies isolate none. Isolation is modeled here as an
in-episode bookkeeping effect with a reward penalty, and, per
Section~\ref{sec:aegis_rl}, has a real containment consequence only on
ToN-IoT-Network -- on CICIoT2023 and UNSW-NB15 no real host identifier
exists in the released data, so "isolating a host" in those two datasets is
a per-flow bookkeeping action with the reward consequence of isolation but
none of its operational one. On ToN-IoT-Network specifically, where
isolation is real, a false isolation would remove a machine from the
network; at a $54\%$ isolation rate this would be operationally severe if
deployed, and the reward penalty alone is not a substitute for a
human-in-the-loop requirement. Since Candidate A -- the method this paper
actually selects and reports results for -- does not use the RL layer at
all, this risk does not apply to the reported system as selected; we
document it here because it is the reason B and C were not selected, and
because any future work that revisits the RL investigation layer inherits
this constraint.

\section{Conclusion}
\label{sec:conclusion}

We began this revision by auditing our prior graph-RL intrusion-detection
study and confirming, with a specific structural diagnosis, that an
RL-as-classifier design cannot remain competitive: an RL policy trained on
the seed/update budgets tractable for one study cannot out-learn a mature
gradient-boosted-tree ensemble at the classification task itself.
\textbf{We therefore redesigned the method as AegisGraph-X, which never
asks RL to classify}: a nine-expert tree/boosting/neural/graph pool supplies
detection, a mixture-of-experts router combines it, and a budgeted RL layer
only decides whether to commit to the calibrated ensemble's own label,
acquire more information, escalate, isolate, or defer -- a structural
guarantee, not a tuning fix, against the failure mode identified in our
audit. Validation-only selection among a backbone-only candidate (A), a
budgeted-RL candidate (B), and a conformal-safety-shielded candidate (C)
selected \textbf{Candidate A}: on CICIoT2023's extreme class imbalance, the
RL layer collapsed to a degenerate always-isolate policy under common reward
weights, and on no dataset did the RL layer improve classification accuracy
over the backbone alone -- we report this honestly rather than search for a
configuration that looks better. \textbf{The selected backbone reaches test macro-F1 $0.940$ (UNSW-NB15),
$0.997$ (ToN-IoT-Network), and $0.949$ (CICIoT2023). Against the strongest
gradient-boosted-tree baseline it is equivalent to within one tenth of a
macro-F1 point on the first two, and reliably worse by $1.2$ points on the
third.} That removes the 20--53 point deficit of the legacy design, which is
a real change. It does not establish that the redesign is better than an
off-the-shelf gradient boosting library, and on CICIoT2023 it is measurably
worse than one.

What the evidence does not support is the more interesting part. Every layer
we added fails its own ablation, and at ten seeds two of those failures are
statistically significant rather than merely nominal. The single best expert
outperforms the full nine-expert pool by $0.0031$ macro-F1
($p = 0.008$, $n = 30$): the ensemble is not simply redundant but actively
worse than choosing one of its members. The learned router is
indistinguishable from an unweighted average ($p = 0.90$) -- a clean null.
The RL investigation layer never spends its telemetry budget on any dataset
under either reward mode, collapses to always-isolate under extreme
imbalance, and improves classification accuracy nowhere. Conformal FPR
control moves alert volume by $0.2\%$. The distilled student's three
hard-class components together cost it $5.6$ macro-F1 points relative to
plain knowledge distillation, though focal loss alone is clearly beneficial.
Taken together, the five-layer architecture is not vindicated by these
experiments; the components that survive are the ones that were already
standard.

We therefore offer this paper as a structural diagnosis and a negative
result rather than as a method. The diagnosis we do defend: RL should not be
asked to own the classification decision in an IDS that must also compete
with mature supervised classifiers, because it will lose that comparison for
reasons that are structural rather than reward-design accidents. The
redesign that follows from the diagnosis is sound in the narrow sense that it
eliminates the failure mode by construction -- and unvalidated in the sense
that none of the machinery it adds on top has yet been shown to earn its
place. We tested the most likely explanation for the investigation layer's
failure directly (Section~\ref{sec:results_analytic}): a calibrated posterior
and an explicit cost table already determine the optimal action
analytically, with no training, and that closed-form policy never isolates
on any dataset and beats Candidate A significantly on two of three. The
learned policy's CICIoT2023 collapse is therefore a training failure, not a
property of what the reward table actually rewards -- a more precise, and
less flattering, conclusion than the uncertainty we reported in earlier
drafts of this section. It does not tell us why PPO fails to find the
optimum the landscape contains, which remains the priority for future work
if the investigation layer is to be revisited at all.

\section*{Statements}

\textbf{Data and code availability.} All three primary corpora are publicly
available~\cite{moustafa2015unswnb15,alsaedi2020toniot,neto2023ciciot2023}.
Experiment code, per-seed result CSVs, split manifests, generated tables, and
run logs accompany this submission; split manifests record row counts, class
balance, and deduplication counts per split for every dataset.

\textbf{Reproducibility caveat.} A network-weight-initialization determinism
defect, disclosed in Section~\ref{sec:results_rewardmodes}, meant that in an
earlier run the nominal seed did not fully control initial conditions. The
defect is fixed; all reported statistical comparisons are computed from runs
executed after the fix, so that the pairing assumption underlying the paired
tests holds.

\textbf{Ethics.} This work uses only public, previously collected network
traces containing no personal data, and involves no human subjects. It
develops a defensive capability. Section~\ref{sec:limitations} states the
deployment risk we consider material: at the ten-seed evaluation scale the
investigation layer isolates a majority of hosts on every dataset, and we do
not consider it deployable without analyst confirmation of isolation
actions. This risk applies only to the (unselected) Candidates B/C; the
selected Candidate A does not isolate hosts at all.

\textbf{Compute.} All experiments ran on a single NVIDIA RTX 5090 (24\,GB)
with a 24-core CPU.

\textbf{AI tool usage.} Large language model assistance was used for
manuscript editing, for the statistical re-analysis reported in
Section~\ref{sec:results_equivalence}, and for code review. All experimental
results, all reported numbers, and all scientific claims were produced and
verified by the authors.

\bibliographystyle{IEEEtran}
\bibliography{refs}

\end{document}
